\documentclass[../../../include/open-logic-section]{subfiles}

\begin{document}

\olfileid{sth}{z}{union}
\olsection{Gabungan}

\olref[sth][z][sep]{prop:intersectionsexist} wis menehi kita
irisan. Nanging supaya gabungan saka sembarang himpunan ana,
kita butuh aksioma liyane:

\begin{axiom}[Gabungan] Kanggo sembarang himpunan $A$, himpunan
$\bigcup A = \Setabs{x}{(\exists b \in A) x \in b}$ ana.
\[
	\forall A \exists U \forall x(x \in U \liff (\exists b \in A)x \in b)
\]
\end{axiom}

Aksioma iki uga didhukung dening konsepsi kumulatif-iteratif.
Upamane $A$ iku himpunan, mula $A$ kawangun ing sawenehing
tataran $S$ (miturut \stageshier). Saben anggota $A$ wis
kawangun \emph{sadurunge} $S$ (miturut \stagesacc); mula
kanthi penalaran kang padha, saben anggota saka saben
anggota $A$ uga kawangun sadurunge $S$. Dadi kabeh
himpunan \emph{iku} wis kasedhiya sadurunge $S$ kanggo
diwangun dadi sawijining himpunan ing $S$. Himpunan iku
persis $\bigcup A$.

\end{document}
